% !TeX root = ../Book1.tex
% 两条同端点路径及其组成的小正向回路。
\begin{figure}[htbp]
	\centering
	\begin{tikzpicture}[font=\small,scale=1.1]
		\fill[GeometryBlueLight] (0,0) rectangle (3,2.1);
		\draw[->,GeometryGray] (-0.4,0)--(3.8,0) node[right] {\(x^1\)};
		\draw[->,GeometryGray] (0,-0.4)--(0,2.8) node[above] {\(x^2\)};
		\draw[GeometryBlue,very thick,->] (0,0)--(1.7,0);
		\draw[GeometryBlue,very thick] (1.7,0)--(3,0);
		\draw[GeometryBlue,very thick,->] (3,0)--(3,1.3);
		\draw[GeometryBlue,very thick] (3,1.3)--(3,2.1);
		\draw[orange!80!black,very thick,->] (0,0)--(0,1.3);
		\draw[orange!80!black,very thick] (0,1.3)--(0,2.1);
		\draw[orange!80!black,very thick,->] (0,2.1)--(1.7,2.1);
		\draw[orange!80!black,very thick] (1.7,2.1)--(3,2.1);
		\fill (0,0) circle[radius=1.8pt] node[below left] {\(p\)};
		\fill (3,2.1) circle[radius=1.8pt] node[above right] {\(q\)};
		\node[below] at (3,0) {\(\varepsilon\)};
		\node[left] at (0,2.1) {\(\varepsilon\)};
		\node[GeometryBlue] at (2.3,0.35) {\(\alpha\)};
		\node[orange!80!black] at (0.65,1.7) {\(\beta\)};
		\node[align=center] at (1.6,1.0) {先沿 \(\alpha\)，再沿 \(\beta^{-1}\)\\\(P_{\beta}^{-1}P_{\alpha}\)};
	\end{tikzpicture}
	\caption{两条路径都从 \(p\) 到 \(q\)。用 \(\beta\) 的逆向运输把 \(\alpha\) 的结果带回 \(E_p\)，得到正向小回路的返回变换。它与恒等变换之差的面积首项为 \(-\varepsilon^2F_{12}(p)\)；两边长度在图中仅作示意。}
	\label{g1:fig:curvature-loop}
\end{figure}
